\section{Introduction}
\label{sec:intro}

Anonymous credentials (AC)~\cite{chaum85-anonymous-credentials,EC:Brands95,EC:CamLys01} allow parties to
convey certified information about themselves in a privacy-preserving way.  The usual roles are issuers,
holders, and verifiers.  A holder obtains attributes certified by an issuer and later proves statements about
those attributes to a verifier without revealing the credential itself.  The central security goals are
unforgeability and unlinkability: verifiers should only accept credentials derived from an issuer, while issuers
and verifiers should not be able to link a presentation to an issuance session or to another presentation.

Certain AC use cases, such as anonymous rate limiting~\cite{ietf-privacypass-arc-protocol-01,schlesinger-cfrg-act-00},
require that a credential be usable only a bounded number of times in a given context, for example per
epoch or per service.  Existing classical constructions can provide compact credentials and presentations, but
practically efficient post-quantum ACs remain much less developed.  This gap is relevant for digital identity
systems, including wallet deployments that will eventually require post-quantum migration.

\paragraph*{Building ACs.}
Anonymous credentials are often built from blind-signature-style issuance.  The holder first commits to the
message to be certified, the issuer signs a target derived from this committed message, and the holder later
presents a zero-knowledge proof of knowledge of a valid message-signature pair.  A generic post-quantum
combination of standard signatures, commitments, and zero-knowledge proofs is currently too large for the
settings targeted here.  We therefore combine a proof-friendly lattice signature, an Ajtai-style commitment, and
SmallWood zero-knowledge proofs in a relation that is designed to be compact.

\subsection{Contributions}
We construct a post-quantum blind-issuance and anonymous rate-limited credential design from lattice
assumptions and hash-based zero knowledge.  The construction builds on proof-friendly lattice signatures
from vanishing SIS with rational hashing~\cite{EPRINT:DKLW25,PKC:DKLW25}, NTRU-style trapdoor sampling
for hash-and-sign signatures~\cite{EPRINT:ENSTW23}, and the SmallWood proof system for relatively small
statements~\cite{EPRINT:FenRiv25b}.  Rate limiting is implemented by signing a hidden PRF key and proving
correct tag derivation during presentation.

The blind-issuance layer uses an Ajtai commitment and a two-party randomness step.  The holder commits
to the full signed message and its holder-side randomness before the issuer sends its contribution.  The final
rational-hash randomness is obtained by centering the sum of the two contributions.  This preserves the shared-
randomness issuance architecture, while the rational hash itself is the BB-tran form of DKLW:
\[
  T=B_0+B_1\vecmu+B_2r_0+Z,\qquad (B_3-r_1)\Had Z=1.
\]
The target is computed from the bare hidden witness elements, not from the commitment value.  The commitment
serves to bind the holder to its message and holder-side randomness before the issuer randomness is known.

The privacy of the public target is obtained by choosing the linear-side randomizer $r_0$ long enough for a
cyclotomic leftover-hash lemma to apply, so that $B_2r_0$ statistically masks the message-dependent and
inverse-side offsets.  The current draft proves correctness and issuer-view hiding of the issuance surface and then
builds a conditional ARC security decomposition from that issuance layer, NIZK knowledge soundness / zero
knowledge, and PRF security.  One-more unforgeability of the interactive issuance layer remains a separate
reduction goal.

\subsection{Technical overview}
The full signed message is
\[
  \vecmu=\vecm\Vert\veck,
\]
where $\vecm$ is the attribute vector and $\veck$ is the secret PRF key.

\paragraph*{Issuance.}
The holder samples $\veck$, holder-side hash randomness $r_0^H,r_1^H$, and commitment-only randomness
$\bar r$.  It sends the Ajtai commitment
\[
  \vecc=\mACom[\vecmu\Vert r_0^H\Vert r_1^H\Vert \bar r]^\top
\]
to the issuer.  The issuer samples bounded randomness $r_0^I,r_1^I$ and returns it.  The holder computes
\[
  r_0=\centerfunc(r_0^H+r_0^I),\qquad
  r_1=\centerfunc(r_1^H+r_1^I),
\]
then derives
\[
  Z=(B_3-r_1)^{-1},\qquad
  T=B_0+B_1\vecmu+B_2r_0+Z.
\]
The holder sends $T$ together with a pre-sign NIZK proving that the same bounded witness opens the
commitment, computes the centered randomness, and satisfies the BB-tran target relation.  If the proof
verifies, the issuer samples a short preimage $\vecu\gets\SampPre(\trapdoor,T)$ such that $\mA\vecu=T$ and
returns it to the holder.

\begin{figure}[h]
  \centering
  \resizebox{\columnwidth}{!}{%
  \begin{tikzpicture}[
    actor/.style={font=\bfseries, align=center},
    lifeline/.style={draw, thick},
    msg/.style={-Latex, thick},
    note/.style={draw, rounded corners, fill=gray!10, align=center, inner sep=4pt},
    x=1.15cm, y=0.9cm]
    \node[actor] (I) at (0,10) {\Large Issuer\\[-1pt]$\Issuer$};
    \node[actor] (H) at (7,10) {\Large Holder\\[-1pt]$\Holder$};
    \draw[lifeline] (I) -- ++(0,-9.5);
    \draw[lifeline] (H) -- ++(0,-9.5);
    \node[note] at ($(H)+(0,-1.3)$) {Sample $\veck$ and set $\vecmu=\vecm\Vert\veck$\\Sample $r_0^H,r_1^H,\bar r$};
    \draw[msg] ($(H)+(0,-2.7)$) -- node[above,sloped] {$\vecc=\mACom[\vecmu\Vert r_0^H\Vert r_1^H\Vert\bar r]^\top$} ($(I)+(0,-2.7)$);
    \node[note] at ($(I)+(0,-3.8)$) {Sample issuer randomness\\$r_0^I,r_1^I$};
    \draw[msg] ($(I)+(0,-4.7)$) -- node[above,sloped] {$r_0^I,r_1^I$} ($(H)+(0,-4.7)$);
    \node[note] at ($(H)+(0,-5.8)$) {Center-combine\\$r_b=\centerfunc(r_b^H+r_b^I)$ for $b\in\{0,1\}$\\$Z=(B_3-r_1)^{-1}$};
    \draw[msg] ($(H)+(0,-7.0)$) -- node[above,sloped] {$T=B_0+B_1\vecmu+B_2r_0+Z,\;\pi_{\mathsf{pre}}$} ($(I)+(0,-7.0)$);
    \node[note] at ($(I)+(0,-8.0)$) {Trapdoor sample\\$\vecu\gets\SampPre(\trapdoor,T)$\\such that $\mA\vecu=T$};
    \draw[msg] ($(I)+(0,-8.9)$) -- node[above,sloped] {Signature preimage $\vecu$} ($(H)+(0,-8.9)$);
    \node[note] at ($(H)+(0,-9.6)$) {Store credential witness\\$(\vecu,\vecm,\veck,r_0,r_1)$};
  \end{tikzpicture}}
  \caption{Issuance flow.  The pre-sign proof binds the BB-tran target to the committed message and the centered shared randomness, while hiding the holder witness.}
  \label{fig:issuance-flow}
\end{figure}

\paragraph*{Showing.}
To present a credential, the holder chooses a nonce, computes
$\prftag=\PRF(\veck,\nonce)$, and proves knowledge of
$\vecu,\vecm,\veck,r_0,r_1,Z$ such that
\begin{align*}
  &(B_3-r_1)\Had Z=1,\\
  &\mA\vecu=B_0+B_1(\vecm\Vert\veck)+B_2r_0+Z,\\
  &\prftag=\PRF(\veck,\nonce).
\end{align*}
The verifier sees only $(\nonce,\prftag,\nizkShow)$ and enforces the rate policy by recording accepted
$(\nonce,\prftag)$ pairs.

\begin{figure}[h]
  \centering
  \resizebox{\columnwidth}{!}{%
  \begin{tikzpicture}[
    actor/.style={font=\bfseries, align=center},
    lifeline/.style={draw, thick},
    msg/.style={-Latex, thick},
    note/.style={draw, rounded corners, fill=gray!10, align=center, inner sep=4pt},
    x=1.15cm, y=0.9cm]
    \node[actor] (H) at (0,0) {\Large Holder\\[-1pt]$\Holder$};
    \node[actor] (V) at (9,0) {\Large Verifier\\[-1pt]$\Verifier$};
    \draw[lifeline] (H) -- ++(0,-7.5);
    \draw[lifeline] (V) -- ++(0,-7.5);
    \node[note] at ($(H)+(0,-2.4)$) {Choose unused $\nonce$\\Compute $\prftag\leftarrow\PRF(\veck,\nonce)$\\Compute $\nizkShow$ proving knowledge of\\$\vecu,\vecm,\veck,r_0,r_1,Z$ satisfying\\$(B_3-r_1)\Had Z=1$\\$\mA\vecu=B_0+B_1(\vecm\Vert\veck)+B_2r_0+Z$\\and $\prftag=\PRF(\veck,\nonce)$};
    \draw[msg] ($(H)+(0,-4.4)$) -- node[above,sloped] {$(\nonce,\prftag,\nizkShow)$} ($(V)+(0,-4.4)$);
    \node[note] at ($(V)+(0,-5.8)$) {Reject if $\nizkShow$ fails\\else reject if $(\nonce,\prftag)$ already seen\\else store $(\nonce,\prftag)$ and accept};
    \draw[msg] ($(V)+(0,-7.0)$) -- node[above,sloped] {accept/reject} ($(H)+(0,-7.0)$);
  \end{tikzpicture}}
  \caption{Showing flow.  The tag is deterministic for a fixed key and nonce, enabling stateful rate limiting, while the proof hides the credential witness.}
  \label{fig:showing-flow}
\end{figure}

\paragraph*{SmallWood proofs.}
We use SmallWood~\cite{EPRINT:FenRiv25b} as the NIZK for the issuance and showing relations.  Both
relations are expressed over committed witness rows packed on an evaluation domain.  The issuance proof
checks the Ajtai opening, the centering equations, and the BB-tran target relation.  The showing proof adds
the short preimage equation and the PRF-tag relation.  The exact compiled statements are stated in
Section~\ref{sec:smallwood-model}; proof-system details are deferred to Appendix~\ref{app:smallwood}.
